\section{\texorpdfstring{Profile $(5,5,5,5)$}{Profile (5,5,5,5)}: local-filter orbits}\label{sec:ppt3-surface}

How little partial-transpose rank can a fully-PPT GME state have? We deform the GHZ-type vector and tune the weights to
make both partial-transpose parity blocks singular. One of the two remaining parameters is an invertible local filter,
so a single parameter controls PPT and GME. Appendix~\ref{sec:ppt3-app-surface} gives the proofs.

The first family has partial-transpose ranks seven or eight. A fully-PPT state of rank at most four is
separable~\cite{KCKL,HLVC,KL,ABLS}, and the same theorem bounds each partial-transpose rank from below: if $\rank\rho^{T_c}\le4$, then $\rho^{T_c}$, and hence $\rho$, is separable across that cut~\cite{KCKL,HLVC,KL}.
Thus $(5,5,5,5)$ is the smallest possible profile for fully-PPT GME. The construction below attains it, including at
extreme points of the fully-PPT set.

For real $\alpha>0$, $0<\beta<\alpha$ and $\gamma>0$ put
$h=\sqrt{\alpha^2+\beta^2}$, $\kappa=\alpha^2-\beta^2$, $s=\alpha/\gamma$ and
\begin{equation}\label{eq:ppt3-weights}
O=\frac \kappa\beta,\qquad E=\frac{\kappa(h+\beta)^2}{\alpha^2\gamma},\qquad G=\frac{\alpha(E-\alpha s)}\beta,\qquad
V=\frac{\alpha^2E^2-2\alpha s\kappa E+(\alpha s)^2\kappa}{\beta\kappa},
\end{equation}
\begin{equation}\label{eq:ppt3-rho-surface}
\rho(\alpha,\beta,\gamma)=\gamma\,u_su_s\dg+\beta\,ww\dg+V\,vv\dg+G\,(wv\dg+vw\dg)+E\,ee\dg+O\,oo\dg,
\end{equation}
with $u_s,v,w,e,o$ as in Section~\ref{sec:ppt3-family}. The coefficients are chosen
to make both partial-transpose blocks singular. Where they come from (at $\alpha=\gamma=1$, so $\beta=\tau$): the odd block of $\rho^{T_A}$ has positive
$2\times2$ leading minor $3\tau O$, so it has rank two only if its $3\times3$ leading minor $3O\,(\tau(\tau+O)-1)$ vanishes, which gives $O=\kappa/\tau$; the last
row and column must then lie in the span of the others, which gives $G=(E-1)/\tau$ and $V=E^2/O+\tau G^2=(E^2-2\kappa E+\kappa)/(\tau\kappa)$ for every $E$, and
the remaining parameter $E$ is fixed by making the even block singular. The
operator is homogeneous of degree one in $(\alpha,\beta,\gamma)$;
its essential parameters are $\tau=\beta/\alpha\in(0,1)$
and $a=\gamma/\alpha>0$. Write
$\rho_{\tau,a}=\rho(1,\tau,a)$ and
$\rho_\tau=\rho_{\tau,1}$, with coefficient matrix $X_\tau$. These symbols denote
\emph{unnormalised} operators, not the normalised state $\rho_t$ and matrix $X_t$ of Section~\ref{sec:ppt3-family}; division by their trace gives physical density matrices. Positive scaling
does not change PPT, rank or separability.

\begin{lemma}[a $\tau$-independent supported witness]\label{lem:ppt3-h0}
\cls{proved analytically; SW-7}
Let $H_0$ be the Hermitian form on $\mathbb C^5$ in the basis $(u,v,w,e,o)$ with diagonal
$(-24-8\sqrt3,\ 1,\ 9,\ 12+4\sqrt3,\ 12+12\sqrt3)$ and $(H_0)_{vw}=(H_0)_{wv}=-\nu$, $\nu=7+4\sqrt3$. Put
$y_1=3+2\sqrt3$, $x_0=2\sqrt3-3$, $p=26+16\sqrt3$, $q=\nu^2$ (so $x_0^2q=y_1^2$), and for $z=x+iy$ let $\xi=x^2$,
$\zeta=y^2$, $r=\xi+\zeta$. Then, with $q_1$ of \eqref{eq:ppt3-curve},
\begin{equation}\label{eq:ppt3-h0}
r\,q_1\dg H_0q_1=\zeta(r-1)^2(r-y_1)^2+\xi(r-x_0)^2(r^2+pr+q)+16\nu\,\xi\zeta r\ \ge0 ,
\end{equation}
so $H_0$ is nonnegative on all cut-product rays of $R$, and $\tr(H_0X)\ge0$ whenever $BXB\dg$ is biseparable.
\end{lemma}
\emph{Idea of the proof.} At fixed $r=|z|^2$ the form $q_1\dg H_0q_1$ is a quadratic in $\xi=x^2$, so it is fixed by its two axis values
($z$ imaginary and $z$ real) and its leading coefficient. Each axis value is a perfect square times a manifestly nonnegative factor.
The computation is in Appendix~\ref{sec:ppt3-app-surface}.

\begin{theorem}[fully-PPT GME states of profile $(5,5,5,5)$]\label{thm:ppt3-orbit}
\cls{proved analytically; FP-1, SW-4, SW-5, SW-6, SW-7}
For $1/\sqrt3<\tau<1$ and every $a>0$, the normalised
$\rho_{\tau,a}$ is fully PPT and GME, with rank profile
$(5,5,5,5)$. The family has the following structure.
\begin{enumerate}[label=(\alph*)]
\item Each operator is cyclically and parity invariant,
with range contained in $R_{1/a}$, and its partial
transposes have rank at most five. In the stated interval
the range equals $R_{1/a}$.
\item For $F_\lambda=\operatorname{diag}(1,\lambda)^{\otimes3}$,
\[
\rho(\alpha,\beta,a\alpha)
=aF_\lambda\,\rho(\alpha,\beta,\alpha)F_\lambda^\dagger,
\qquad \lambda^2=a^{-1}.
\]
Thus varying $a$ is an invertible local filter;
PPT and GME depend only on $\tau$.
\item The window is sharp within this ansatz:
$\rho_{\tau,a}$ is not positive for $0<\tau<1/\sqrt3$.
At $\tau=1/\sqrt3$ it has rank four and is separable
across every cut. As $\tau\uparrow1$ its normalisation
tends to a fully separable state
(Remark~\ref{rem:ppt3-endpoint}).
\end{enumerate}
The GME proof uses the fixed supported witness $H_0$
of Lemma~\ref{lem:ppt3-h0} when $1/\sqrt3<\tau\le3/4$, and
$H_5$ from Proposition~\ref{prop:ppt3-witness}
when $3/4\le\tau<1$; the case $a\ne1$ follows because the local filter of (b)
preserves biseparability, full PPT and ranks.
\end{theorem}

\emph{Idea of the proof.} Homogeneity and the local filter reduce the problem to $\alpha=\gamma=1$. With
$Q=\tau+\sqrt{1+\tau^2}$, the odd parity block of $\rho_\tau^{T_A}$ has an explicit Gram factorisation of rank two, and the even block
has vanishing determinant and a third leading minor that changes sign at $\tau=1/\sqrt3$. Together with positivity of
$\rho_\tau$, these give the sharp window and ranks; cyclic symmetry covers the other cuts. The two witnesses cover
the window by one-variable sign checks.

\begin{corollary}[an explicit extreme state with rational weights and profile $(5,5,5,5)$]\label{cor:ppt3-extreme5}
\cls{exact-rational certificate, computer-assisted; the witness proof of genuine entanglement is exact-symbolic; extremality on the whole surface is proved from the $2\times4$ low-rank theorem and independently certified by exact-symbolic block minors; FP-2, FP-3, FP-4, FP-5, SW-3, DS-6, SB-1, SB-2}
Let $\tilde u=2\ket{000}-\ket{011}-\ket{101}-\ket{110}$, let $v,w,e,o$ be as in \eqref{eq:ppt3-vectors}, and
\begin{multline}\label{eq:ppt3-sigma5}
N=825000\,\tilde u\tilde u\dg+1440000\,ww\dg+2021184\,vv\dg\\
+145200\,(wv\dg+vw\dg)+951720\,ee\dg+450625\,oo\dg,
\qquad\sigma=\frac N{16323219}.
\end{multline}
Then $\sigma$ has rank five and is invariant under cyclic permutations and global parity; $N=30000\,\rho(55,48,110)$ is
the member $(\alpha,\beta,\gamma)=(55,48,110)$ of Theorem~\ref{thm:ppt3-orbit}, the Pythagorean choice $(m,n,a)=(8,3,2)$ with
$\alpha=m^2-n^2$, $\beta=2mn$ and $\gamma=a\alpha$.
Here $\tau=48/55$ and $s=1/2$; all six coefficients
in \eqref{eq:ppt3-rho-surface} are rational.
The projector coefficients are rational, but the matrix entries of $\sigma$ in the computational basis
belong to $\mathbb Q(\omega)$, since $e$ and $o$ contain cube-root phases.
All three partial transposes of $\sigma$ are positive semidefinite of rank five (at $A|BC$ the even and odd
blocks have ranks $3$ and $2$). The state $\sigma$ is an extreme point of the set of fully-PPT three-qubit states
and is GME, and $(1-p)\sigma+pI/8$ is GME for every $p<4.05\times10^{-3}$ (the certificate is in Appendix~\ref{sec:ppt3-app-surface}, using
the witness of Corollary~\ref{cor:ppt3-noise-reused}). More generally, \emph{every} normalised $\rho_{\tau,a}$ of Theorem~\ref{thm:ppt3-orbit}
($1/\sqrt3<\tau<1$, $a>0$) is an extreme point of the set of fully-PPT three-qubit states, and an edge state across every cut: no product
vector across a cut lies in its range with its partial conjugate in the range of the corresponding partial transpose.
\end{corollary}

\emph{Idea of the proof.} Theorem~\ref{thm:ppt3-orbit} gives positivity and ranks, and $H_{1/2}$ detects GME directly.
The noise bound uses the level-two witness $Y_*$ of Corollary~\ref{cor:ppt3-noise-reused}. For extremality, a $2\times4$ PPT state of birank $(5,5)$ that is entangled across a cut is extreme in that cut's PPT set, because a state on the boundary of its face has rank or partial-transpose rank at most four
and is then separable~\cite{KCKL,HLVC}; independently, the cyclic, parity and
reflection symmetries split the rank-$24$ face-direction calculation into blocks whose minors are explicit polynomials in
$Q=\tau+\sqrt{1+\tau^2}$ with no root in the window (Appendix~\ref{sec:ppt3-app-surface}).

\begin{proposition}[the profile-$(5,5,5,5)$ states are not isolated]\label{prop:ppt3-moduli}
\cls{exact-symbolic certificate, computer-assisted; the smoothness and openness steps are proved from the exact ranks; MO-1, MO-2, MO-3, MO-4, SB-1, SB-3, SB-4}
Near every state $\rho$ of Theorem~\ref{thm:ppt3-orbit} (in particular near $\sigma$ of Corollary~\ref{cor:ppt3-extreme5}), the fully-PPT three-qubit states of rank profile $(5,5,5,5)$ form
a smooth real manifold $\mathcal M$ of dimension $27$, and every state of $\mathcal M$ near $\rho$ is GME and an extreme
point of the set of fully-PPT states. Invertible local filters ($\rho\mapsto F\rho F\dg/\tr(F\rho F\dg)$,
$F=F_1\otimes F_2\otimes F_3$) act on $\mathcal M$ with orbits of dimension $18$, so $\mathcal M$ has nine local
moduli. The construction of Theorem~\ref{thm:ppt3-orbit} is only a two-parameter
cyclic-symmetric family (one parameter modulo filtering), not a classification.
\end{proposition}

\emph{Idea of the proof.} Rank-five states have $54$ real parameters. Each rank-five partial transpose imposes
$(8-5)^2=9$ smooth equations; the symmetry blocks of their differentials have explicit minors in $Q$ that do not vanish on the window, so they are
independent at every state of the surface. The implicit function theorem gives dimension $27$, and local filters remove
$18$ parameters (the stabiliser of the support is finite, Corollary~\ref{cor:ppt3-chirality}). GME and extremality persist nearby
(Appendix~\ref{sec:ppt3-app-surface}).

\begin{remark}[the two endpoints; two witnesses are needed]\label{rem:ppt3-endpoint}
\cls{proved; endpoints and uniqueness checked symbolically in SW-7, SW-8}
At $\tau=1/\sqrt3$ ($Q=\sqrt3$) the operator $\rho_\tau$ is a limit of fully-PPT states, hence positive semidefinite with all three
partial transposes positive semidefinite, and it has rank four, so it is separable across each cut by the $2\times4$
low-rank theorem~\cite{KCKL,HLVC}. Explicitly, at $a=1$ its coefficient matrix
$X_{1/\sqrt3}=\operatorname{diag}(1,3\sqrt3,\sqrt3/3,2,2\sqrt3/3)+\sqrt3\,(E_{vw}+E_{wv})$ is a positive combination of four
$A|BC$-product rays, at the zeros $z=\pm i\sqrt{y_1},\pm\sqrt{x_0}$ of $q_1(z)\dg H_0q_1(z)$ in~\eqref{eq:ppt3-h0}:
\begin{equation}\label{eq:ppt3-lower-endpoint}
 X_{1/\sqrt3}
  =\frac{3\sqrt3-5}{48}\sum_{\pm}q_1\bigl(\pm i\sqrt{y_1}\bigr)q_1\bigl(\pm i\sqrt{y_1}\bigr)\dg
  +\frac{3\sqrt3+5}{48}\sum_{\pm}q_1\bigl(\pm\sqrt{x_0}\bigr)q_1\bigl(\pm\sqrt{x_0}\bigr)\dg ,
\end{equation}
with weights $(3\sqrt3\mp5)/48>0$. Cyclic permutation gives the decompositions across $B|CA$ and $C|AB$, and the local filter in
Theorem~\ref{thm:ppt3-orbit}(b) carries them to every $a>0$.
It is not fully separable. Writing the product vectors of Proposition~\ref{prop:ppt3-witness}(a) as
$B_sq_s(z)=(\ket0+z\ket1)\otimes\phi_s(z)$, the $BC$ coefficient matrix of $\phi_s(z)$ has determinant
$-3(z^2+s)^2$, so the only fully product rays in $R_s$ are $\ket{111}$ and $(\ket0\pm i\sqrt s\ket1)^{\otimes3}$.
A fully separable state with range in $R_{1/a}$ therefore has rank at most three, while $\rho_{1/\sqrt3,a}$ has rank four.

At the upper endpoint, for every fixed $a>0$ the limit is fully separable and has rank profile $(2,2,2,2)$.
Indeed the weights satisfy $(O,E,G,V)\to(0,0,-1/a,1/a^2)$ as $\tau\uparrow1$, giving
\[
 \rho_{1,a}:=\lim_{\tau\uparrow1}\rho_{\tau,a}
    =a\,u_{1/a}u_{1/a}\dg+(w-a^{-1}v)(w-a^{-1}v)\dg.
\]
Put $\ket{\eta_\pm(a)}=(\sqrt a\ket0\pm i\ket1)/\sqrt{a+1}$. Expanding the product vectors gives
\begin{equation}\label{eq:ppt3-upper-endpoint}
 \frac{\rho_{1,a}}{\tr \rho_{1,a}}
   =\frac12\left[
     (\ket{\eta_+(a)}\bra{\eta_+(a)})^{\otimes3}
    +(\ket{\eta_-(a)}\bra{\eta_-(a)})^{\otimes3}\right],
 \qquad \tr \rho_{1,a}=\frac{(a+1)^3}{a^2}.
\end{equation}
The two local rays are distinct for every finite $a>0$, and remain distinct after conjugating any subset of
factors. Thus both the state and its three partial transposes have rank exactly two. The GME interval therefore
terminates at states separable across every cut at both ends; only the upper endpoint is fully separable.

Together, \eqref{eq:ppt3-lower-endpoint} and~\eqref{eq:ppt3-upper-endpoint} explain the origin of $H_0$ and show that no
single witness detects $\rho_\tau$ on the whole window $1/\sqrt3<\tau<1$, so two witnesses are needed. Only the compression of a
witness to $R$ matters, and it is a form $H$ on $\mathbb C^5$ nonnegative on the cut-product rays of $R$. If $\tr(HX_\tau)<0$
on the whole window, continuity gives $\tr(HX_{1/\sqrt3})\le0$ and $\tr(HX_1)\le0$, where $X_1$ is the coefficient matrix of
$\rho_{1,1}$. Both endpoint matrices are positive
combinations of cut-product rays, so $q_1(z)\dg Hq_1(z)$ vanishes, and being nonnegative has a double zero, at
$z=\pm i\sqrt{y_1},\pm\sqrt{x_0}$ and at $z=\pm i$. The cyclic, parity and $K\circ S_{BC}$ symmetries fix every $X_\tau$ and
permute the cut-product rays, so the average $\bar H$ of $H$ over them has the same properties. These double zeros determine
$\bar H$ up to a nonnegative multiple of $H_0$ (Appendix~\ref{sec:ppt3-app-surface}). But $\tr(H_0X_\tau)$ changes sign at
$\tau\approx0.805$, the root of $P(Q)$ (for instance $\tr(H_0X_{48/55})>0$), a contradiction.
\end{remark}

\begin{remark}[the edge condition fails exactly at the endpoint rays]\label{rem:ppt3-edge-overlap}
\cls{exact-symbolic certificate, computer-assisted; EO-1, EO-2, EO-3, EO-4, EO-5}
Take $a=1$; invertible local filters map product vectors to product vectors and commute with partial conjugation up to the conjugate
filter, so they preserve the edge property. By Proposition~\ref{prop:ppt3-witness}(a) the $A|BC$ product rays of $R_1$ are those of
$\psi(z)=Bq_1(z)=(\ket0+z\ket1)\otimes\phi_1(z)$, $z\in\mathbb C$, and of $v=\ket{111}$. So $\rho_\tau$ is an edge state across $A|BC$ exactly when none of them has its partial conjugate
$\psi^{T_A}(z)=(\ket0+\bar z\ket1)\otimes\phi_1(z)$ in $\operatorname{ran}\rho_\tau^{T_A}$, that is, orthogonal to the kernel vectors $k_1,k_2,k_3$ of Appendix~\ref{sec:ppt3-app-surface}.
For $z=x+iy$,
\[
k_1\dg\psi^{T_A}(z)=6\bigl((Q^2-1)|z|^2+Q(z^2-1)\bigr),\qquad\operatorname{Im}=12Qxy .
\]
A zero therefore has $y=0$, $x^2=Q/(Q^2+Q-1)$, or $x=0$, $y^2=Q/(Q^2-Q-1)$, and there
\[
k_3\dg\psi^{T_A}=-\frac{6Qx(Q^2-3)(Q^2+2Q-1)}{Q^2+Q-1}\quad\text{resp.}\quad-\frac{6iQy(Q^2-3)(Q^2-2Q-1)}{Q^2-Q-1},
\]
while $k_2\dg\ket{111}=12$. On $\sqrt3<Q<1+\sqrt2$ no factor vanishes, so the edge condition holds; this checks the edge property across $A|BC$ directly, and across the other cuts by
the cyclic shift. At the endpoints the factors $Q^2-3$ and $Q^2-2Q-1$ vanish. At $Q=\sqrt3$ the two branches give $z=\pm\sqrt{x_0},\pm i\sqrt{y_1}$, the four rays of~\eqref{eq:ppt3-lower-endpoint}, and at $Q=1+\sqrt2$
the imaginary branch gives $z=\pm i$, the rays of~\eqref{eq:ppt3-upper-endpoint}. So at both ends the common zeros of the overlaps with $k_1,k_2,k_3$ are exactly the product rays of the separable decomposition of the limit state.
\end{remark}

